\documentclass{ltjsbook}
\usepackage{docmute} 
\usepackage{../../myStyle} 

\begin{document}
\chapter{行列と方程式}
\label{L03_Equation}

\section{行列と方程式}
線形代数は「便利な書き方」の学問です。便利な書き方をするためにルールが作られていますから，ルールから学ぶと「なんでそんな変な操作をするんだ」という気持ちになるのもわかります。

では何が便利になるのでしょうか。これは方程式を解くことと関係があります。
たとえば，以下のような連立方程式があったとしましょう。

\[
    \left \{
    \begin{array}{lcl}
        x-2y-5z  & = & 3 \\
        5x+4y+3z & = & 1 \\
        3x+y-3z  & = & 6
    \end{array}
    \right .
\]
これは行列で表現すると，次のようになります。
\[
    \begin{pmatrix} 1 & -2 & -5 \\
                5 & 4  & 3  \\
                3 & 1  & -3\end{pmatrix}
    \begin{pmatrix} x \\ y \\ z \end{pmatrix}
    =
    \begin{pmatrix} 3 \\ 1 \\ 6 \end{pmatrix}
\]

この左辺を行列とベクトルの式の計算ルールにのっとって展開してみてください。ちゃんと最初の連立方程式の左辺になることがわかると思います。かけて足して，という面倒な計算ルールは，連立方程式を簡単に表記するためのものだったのですね。

最終的にはこの方程式を解いて，次のように答えを求めます。
\[
    \begin{pmatrix} 1 & -2 & -5 \\
                5 & 4  & 3  \\
                3 & 1  & -3\end{pmatrix}
    \begin{pmatrix} -1 \\ 3 \\ -2 \end{pmatrix}
    =
    \begin{pmatrix} 3 \\ 1 \\ 6 \end{pmatrix}
\]

\subsection{連立方程式を解こう}
さて，ではこの方程式の解はどのようにして求めたのでしょうか。改めてその計算方法を見てみましょう。解を導く方法の一つ，消去法は，ひとつの方程式を何倍かして，他の方程式に加えることにより，方程式をどんどん簡単にしていくというものです。
まず，第一の式を5倍，あるいは3倍して，第二，第三の式からxの項を消去します。
\[ \left \{
    \begin{array}{rcl}
        x-2y-5z  & = & 3  \\
        -14y-28z & = & 14 \\
        -7y-12z  & = & 3
    \end{array}
    \right .
\]
第二の式の係数を簡単にしておきましょう。
\[ \left \{
    \begin{array}{rcl}
        x-2y-5z & = & 3  \\
        y+2z    & = & -1 \\
        -7y-12z & = & 3
    \end{array}
    \right .
\]
第二の式を7倍して，第三の式からyを消去します。
\[ \left \{
    \begin{array}{rcl}
        x-2y-5z & = & 3  \\
        y+2z    & = & -1 \\
        2z      & = & -4
    \end{array}
    \right .
\]
あとはこれの3行目から$z=-2$が得られ，芋づる式に$x=-1$，$y=3$が得られました。


この操作は，式を一本ずつ，あるいは2つの式を組み合わせて文字を消していく消去法を係数全体に行う操作になっています。実際，ここで操作される係数だけ見ていくと，次のようになっているのでした。

\paragraph{第一段階}；一行目の式を使って，二行目以降の一列目の数字をゼロにする。
\[
    \begin{pmatrix} 1 & -2 & -5  \\
                0 & 1  & 2   \\
                0 & -7 & -12\end{pmatrix}
    \begin{pmatrix} x \\ y \\ z \end{pmatrix}
    =
    \begin{pmatrix} 3 \\ -1 \\ 3 \end{pmatrix}
\]

\paragraph{第二段階}；二行目の数字を使って，三行目二列目の数字をゼロにする。

\[
    \begin{pmatrix} 1 & -2 & -5 \\
                0 & 1  & 2  \\
                0 & 0  & 2\end{pmatrix}
    \begin{pmatrix} x \\ y \\ z \end{pmatrix}
    =
    \begin{pmatrix} 3 \\ -1 \\ -4 \end{pmatrix}
\]

\paragraph{第三段階}；三行目の関係から，対角要素を1にする。

\[
    \begin{pmatrix} 1 & -2 & -5 \\
                0 & 1  & 2  \\
                0 & 0  & 1\end{pmatrix}
    \begin{pmatrix} x \\ y \\ z \end{pmatrix}
    =
    \begin{pmatrix} 3 \\ -1 \\ -2 \end{pmatrix}
\]

\paragraph{第四段階}；三行目の関係を使って二行三列目の数字をゼロにする。

\[
    \begin{pmatrix} 1 & -2 & -5 \\
                0 & 1  & 0  \\
                0 & 0  & 1\end{pmatrix}
    \begin{pmatrix} x \\ y \\ z \end{pmatrix}
    =
    \begin{pmatrix} 3 \\  3 \\ -2 \end{pmatrix}
\]

\paragraph{第五段階}；二，三行目の関係を使って，一行二，三列目の数字をゼロにする。

\[
    \begin{pmatrix} 1 & 0 & 0 \\
                0 & 1 & 0 \\
                0 & 0 & 1\end{pmatrix}
    \begin{pmatrix} x \\ y \\ z \end{pmatrix}
    =
    \begin{pmatrix} -1 \\  3 \\ -2 \end{pmatrix}
\]

最後の形を見ればお分かりと思いますが，連立方程式を解くこととは，\underline{係数の行列を単位行列に置き換える}ことと同じなのです!

この一連の手続きを，行列の言葉で表現し直してみたいと思います。
\subsection{行の基本変形}
行列の行に関して，次の操作を\keytermL{行基本変形}{ぎょうきほんへんけい}と言います。
\begin{enumerate}
    \item ある行に0ではない数字を掛ける。
    \item ある行に，他の行を何倍かしたものを加える。
    \item 二つの行を入れ替える。
\end{enumerate}
連立方程式を解くプロセスは，この3つの操作を駆使して行われてきたことは明らかです。これを行列の計算で表せるようにしましょう。

\paragraph{ある行に0ではない数字を掛ける}
ある行列\bm{A}の$i$行目を$c$倍したいとします。この操作を$\bm{P}_i(c)\bm{A}$とすると，$\bm{P}_i(c)$は次のように表せます。

\[
    \bm{P}_i(c) = \begin{pmatrix}
        1 &        &   &   &   &        &   \\
          & \ddots &   &   &   & \bm{O} &   \\
          &        & 1 &   &   &        &   \\
          &        &   & c &   &        &   \\
          &        &   &   & 1 &        &   \\
          & \bm{O} &   &   &   & \ddots &   \\
          &        &   &   &   &        & 1 \\
    \end{pmatrix}
\]
単位行列の$i$列目だけ$c$が入っているような行列ですね。数値例で確認しておきます。
\[ \bm{A} = \begin{pmatrix} 1 & 2 & 4 & 3 \\ 2 & 1 & -1 & 3 \\ 3 & 1 & 4 & 1 \end{pmatrix} \]
の3行目を2倍する場合，
\[
    \bm{P}_3(2)\bm{A} =
    \begin{pmatrix} 1 & 0 & 0 \\ 0 & 1 & 0 \\ 0 & 0 & 2 \end{pmatrix}
    \begin{pmatrix} 1 & 2 & 4 & 3 \\ 2 & 1 & -1 & 3 \\ 3 & 1 & 4 & 1 \end{pmatrix}
    =
    \begin{pmatrix} 1 & 2 & 4 & 3 \\ 2 & 1 & -1 & 3 \\ 6 & 2 & 8 & 2 \end{pmatrix}
\]
\paragraph{ある行に，他の行を何倍かしたものを加える}
ある行列\bm{A}の$i$行目に，$j$行目の$c$倍したものを加えたいと思います。この操作を表す行列$\bm{P}_{ij}(c)$ は，次のように表せます。
\[
    \bm{P}_{ij}(c) = \begin{pmatrix}
        1 &        &   &        &        &        &   \\
          & \ddots &   &        &        & \bm{O} &   \\
          &        & 1 & \cdots & c      &        &   \\
          &        &   & \ddots & \vdots &        &   \\
          &        &   &        & 1      &        &   \\
          & \bm{O} &   &        &        & \ddots &   \\
          &        &   &        &        &        & 1 \\
    \end{pmatrix}
\]
数値例で確認しておきます。先ほどの行列$\bm{A}$の3行目を2倍して1行目に加えるとしましょう。
\[
    \bm{P}_{13}(2)\bm{A} =
    \begin{pmatrix} 1 & 0 & 2 \\ 0 & 1 & 0 \\ 0 & 0 & 1 \end{pmatrix}
    \begin{pmatrix} 1 & 2 & 4 & 3 \\ 2 & 1 & -1 & 3 \\ 3 & 1 & 4 & 1 \end{pmatrix}
    =
    \begin{pmatrix} 7 & 4 & 12 & 5 \\ 2 & 1 & -1 & 3 \\ 6 & 2 & 8 & 2 \end{pmatrix}
\]

\paragraph{二つの行を入れ替える}
ある行列の$i$行目と$j$行目を入れ替える操作を考えます。これは行列が連立方程式を表していると思えば，連立して成立しているのですから式の順番を入れ替えることに本質的な変化をもたらさないことから正当化しうるでしょう。この操作を行列で$\bm{P}_{ij}$とするなら，次のように表せます。
\[
    \bm{P}_{ij} =
    \begin{pmatrix}
        1 &        &        &        &        &        &        &        &   \\
          & \ddots &        &        &        &        &        & \bm{O} &   \\
          &        & 0      & \cdots & \cdots & \cdots & 1      &        &   \\
          &        & \vdots & 1      &        &        & \vdots &        &   \\
          &        & \vdots &        & \ddots &        & \vdots &        &   \\
          &        & \vdots &        &        & 1      & \vdots &        &   \\
          &        & 1      & \cdots & \cdots & \cdots & 0      &        &   \\
          & \bm{O} &        &        &        &        &        & \ddots &   \\
          &        &        &        &        &        &        &        & 1 \\
    \end{pmatrix}
\]
これも数値例で確認しておきます。先ほどの行列$\bm{A}$の1行目と3行目を入れ替えるとしましょう。
\[
    \bm{P}_{13}\bm{A} =
    \begin{pmatrix} 0 & 0 & 1 \\ 0 & 1 & 0 \\ 1 & 0 & 0 \end{pmatrix}
    \begin{pmatrix} 1 & 2 & 4 & 3 \\ 2 & 1 & -1 & 3 \\ 3 & 1 & 4 & 1 \end{pmatrix}
    =
    \begin{pmatrix} 3 & 1 & 4 & 1 \\ 2 & 1 & -1 & 3 \\ 1 & 2 & 4 & 3 \end{pmatrix}
\]

この通り，連立方程式を解くための操作もすべて行列演算で表現できることがわかりました。
この3つの操作をする行列，$\bm{P}_i(c), \bm{P}_{ij}(c),\bm{P}_{ij}$のことを，\keytermL{基本行列}{きほんぎょうれつ}といいます。
\subsection{連立方程式を基本行列で解こう}
このことを踏まえて，連立方程式の解き方をもう一度くわしく見てみたいと思います。
連立方程式，
\[
    \begin{pmatrix} 1 & -2 & -5 \\
                5 & 4  & 3  \\
                3 & 1  & -3\end{pmatrix}
    \begin{pmatrix} x \\ y \\ z \end{pmatrix}
    =
    \begin{pmatrix} 3 \\ 1 \\ 6 \end{pmatrix}
\]
を行列で表し，
\[ \bm{Ax}=\bm{b} \]
としましょう。ここで\bm{A}は\keytermL{係数行列}{けいすうぎょうれつ}と呼ばれます。
また，計算に際しては\bm{A}に右辺のベクトル\bm{b}を付け加えた行列，すなわち
\[
    [ \bm{A},\bm{b} ] =
    \begin{pmatrix} 1 & -2 & -5 & 3\\ 5 & 4 & 3 & 1 \\3 & 1 & -3 & 6 \end{pmatrix}
\]
を考え，一気に操作してしまいましょう。この行列$ [ \bm{A},\bm{b} ] $は特に連立一次方程式の\keytermL{拡大係数行列}{かくだいけいすうぎょうれつ}と呼ばれます。

さて，今から方程式を解くための操作ステップを全て基本行列で表し，操作のログを取っていくことにします。最初の例と同じ数字ですが，ゴールが係数行列を単位行列にすることだとわかっているので，途中のステップを少しスピードアップしています。

\paragraph{第一段階}；一行目の式を使って，一列目の非対角要素をゼロにする。

\begin{center}
    \begin{tabular}{ll}
        $
            \begin{pmatrix}
                1 & -2 & -5 & 3   \\
                0 & 14 & 28 & -14 \\
                0 & 7  & 12 & -3
            \end{pmatrix}
        $
         &
        \begin{tabular}{l}
            1.\bm{Ab}の1行目を-5倍して2行目に加える，すなわち$\bm{P}_{21}(-5)$ \\
            2.\bm{Ab}の1行目を-3倍して3行目に加える，すなわち$\bm{P}_{31}(-3)$
        \end{tabular} \\
    \end{tabular}
\end{center}

\paragraph{第二段階}；二行目の数字を使って，二列目の非対角要素をゼロにする。
\begin{center}
    \begin{tabular}{ll}
        $
            \begin{pmatrix}
                1 & 0 & -1 & 1  \\
                0 & 1 & 2  & -1 \\
                0 & 7 & 12 & -3
            \end{pmatrix}
        $
         &
        \begin{tabular}{l}
            3.\bm{Ab}の2行目を1/14倍する。すなわち$\bm{P}_{2}(\frac{1}{14})$ \\
            4.\bm{Ab}の2行目を2倍して1行目に加える。すなわち$\bm{P}_{12}(2)$       \\
            5.\bm{Ab}の2行目を-7倍して3行目に加える。すなわち$\bm{P}_{32}(-7)$     \\
        \end{tabular}
    \end{tabular}
\end{center}

\paragraph{第三段階}；三行目の数字を使って，三列目の非対角要素をゼロにする。
\begin{center}
    \begin{tabular}{ll}
        $
            \begin{pmatrix}
                1 & 0 & 0 & -1 \\
                0 & 1 & 0 & 3  \\
                0 & 0 & 1 & -2
            \end{pmatrix}
        $
         &
        \begin{tabular}{l}
            6.\bm{Ab}の3行目を-1/2倍する。すなわち$\bm{P}_{3}(-\frac{1}{2})$ \\
            7.\bm{Ab}の3行目を1倍して1行目に加える。すなわち$\bm{P}_{13}(1)$       \\
            8.\bm{Ab}の3行目を-2倍して2行目に加える。すなわち$\bm{P}_{23}(-2)$     \\
        \end{tabular}
    \end{tabular}
\end{center}

基本行列による操作は8ステップありました。全てを網羅的につなげて書くと，
\[ \bm{P}_{32}(-2)\bm{P}_{31}(1)\bm{P}_{3}(1/2)\bm{P}_{23}(-7)\bm{P}_{21}(2)\bm{P}_{2}(1/14)\bm{P}_{13}(-3)\bm{P}_{12}(-5)\bm{A} = \bm{I} \]
ということになります。係数行列$\bm{A}$の左から右に連なっている点に注意してください。

さてこの操作，長いので
\[ \bm{P} = \bm{P}_{32}(-2)\bm{P}_{31}(1)\bm{P}_{3}(1/2)\bm{P}_{23}(-7)\bm{P}_{21}(2)\bm{P}_{2}(1/14)\bm{P}_{13}(-3)\bm{P}_{12}(-5) \]
とまとめて表現してしまいますが，この操作ログだけに注目してみたいと思います。操作ログだけ取り出すために，単位行列$\bm{I}$に操作ログを適用した$\bm{B} = \bm{PI}$を考えると，次の関係が成り立ちます。
\[ \bm{PIA} =\bm{BA} =  \bm{I} \]
この\bm{B}は，係数行列を単位行列に変えるものであり，いわば行列における逆数のような操作です。

行列でない世界であれば，ある数字$a$の逆数を$a^{-1}$と書きます。ここで
\[ a a^{-1}=a^{-1}a = 1 \]
です。同様に行列の場合にも，正方行列$\bm{A}$に対して
\[ \bm{XA} = \bm{I} \text{かつ} \bm{AX} = I \]
な行列があれば，これを特に\keytermL{逆行列}{ぎゃくぎょうれつ}と呼びます。今回得られた$\bm{B}= \bm{PI}$はまさに$\bm{A}$の逆行列なのです。$\bm{A}$の逆行列は$\bm{A}^{-1}$と表記します。

方程式を解くことは，係数行列を単位行列にすることだと言いました。すなわち，
\[ \bm{Ax} = \bm{b} \]
であるときに，
\[ \bm{Ix} = \bm{b}^{\ast} \]
となれば方程式が解けたといえるわけです。
ここで$\bm{AA}^{-1} = \bm{A}^{-1}\bm{A} = \bm{I}$ですから，
\[ \bm{A}^{-1}\bm{Ax} = \bm{A}^{-1}\bm{b} \]
とすれば方程式が解けたといえることになります。方程式を解くことは，逆行列を求めることだったのです。

\end{document}


